\section{Conclusion}

We present \emph{NetDream}, a graph world-modeling approach that turns Kubernetes autoscaling from reactive threshold control into predictive safe planning. Rather than scaling after pressure appears or buying safety through static over-provisioning, NetDream learns how load, latency, replicas, and failures evolve across a microservice graph, imagines the multi-step consequences of candidate scaling actions, and rejects futures predicted to violate service-level constraints before acting. On a production-grade AWS Kubernetes deployment of Online Boutique, NetDream-Safe achieves the most cost-efficient safe operating point, reducing violations far more efficiently than static replica floors and avoiding the over-provisioning that model-free RL adopts to remain safe. The gain comes from the method's two ingredients: topology-aware world modeling, which makes short-horizon imagination accurate, and explicit safety filtering, which turns prediction into reliable action selection. The takeaway is that safer, cheaper cloud control comes from planning in a learned graph model, not from faster reaction or larger replica floors.